% Exported from Math++. Compile with pdflatex.
\documentclass[11pt,a4paper]{article}
\usepackage[T1]{fontenc}
\usepackage[utf8]{inputenc}
\usepackage{lmodern}
\usepackage{amsmath,amssymb,amsthm,mathtools}
\usepackage{graphicx,booktabs,xcolor}
\usepackage[hidelinks]{hyperref}
\usepackage[margin=25mm]{geometry}
\providecommand{\degree}{\ensuremath{^\circ}}
\providecommand{\xmark}{\ensuremath{\times}}
\providecommand{\omicron}{o}
\providecommand{\oiint}{\mathop{\oint\mkern-12mu\oint}}
\theoremstyle{plain}
\newtheorem{theorem}{Theorem}
\newtheorem{lemma}[theorem]{Lemma}
\newtheorem{corollary}[theorem]{Corollary}
\newtheorem{proposition}[theorem]{Proposition}
\newtheorem{conjecture}[theorem]{Conjecture}
\theoremstyle{definition}
\newtheorem{definition}[theorem]{Definition}
\newtheorem{example}[theorem]{Example}
\theoremstyle{remark}
\newtheorem*{remark}{Remark}
\setlength{\parindent}{0pt}\setlength{\parskip}{0.6em}
\title{Curve sketching}
\date{}
\begin{document}
\maketitle
The analysis helpers answer the questions of a curve-sketching exercise. Each one has a
sentence form and an operator form that mean the same thing.

\section*{A cubic}

\begin{equation*}
  f(x) \coloneqq  x^3 - 3x^2 + 1
\end{equation*}
Find the domain of $f$.

Find the intercepts of $f$.

Find the critical points of $f$.

Find the local extrema of $f$.

Find the intervals of increase of $f$.

Find the concavity of $f$.

Find the absolute extrema of $f$ on $[-1/2, 4]$.

The same questions in operator form:

\begin{equation*}
  \operatorname{Critical\ points} f = \left\{0, 2\right\}
\end{equation*}
\begin{align*}
  \operatorname{Extrema}_{[-1/2, 4]} f = \begin{array}{ll} \text{local maxima} & \left\{\left(0, 1\right)\right\} \\ \text{local minima} & \left\{\left(2, -3\right)\right\} \\ \text{absolute maximum} & \left\{\left(4, 17\right)\right\} \\ \text{absolute minimum} & \left\{\left(2, -3\right)\right\} \end{array}
\end{align*}
\section*{Lines at a point}

Find the tangent line of $f$ at $x = 1$.

Find the normal line of $f$ at $x = 3$.

\begin{equation*}
  \operatorname{Linearisation}_{x=3} f \to L\left(x\right) = 9 x - 26
\end{equation*}
\section*{A rational function}

\begin{equation*}
  g(x) \coloneqq  \frac{2x^2 + 1}{x^2 - 1}
\end{equation*}
Find the domain of $g$.

Find the range of $g$.

Find the asymptotes of $g$.

\begin{align*}
  \operatorname{Asymptotes} \frac{x^2 + 1}{x - 1} = \begin{array}{ll} \text{vertical} & \left\{x = 1\right\} \\ \text{horizontal} & \text{none} \\ \text{slant} & \left\{y = x + 1\right\} \end{array}
\end{align*}
\section*{The mean value and intermediate value theorems}

\begin{align*}
  \operatorname{Mean\ value\ theorem}_{[0, 2]} (x^3 - x) = \begin{array}{ll} \text{continuous on [0, 2]} & \text{yes} \\ \text{differentiable on (0, 2)} & \text{yes} \\ \text{average rate of change} & 3 \\ \text{c with f\prime (c) equal to it} & \left\{\frac{2 \sqrt{3}}{3}\right\} \end{array}
\end{align*}
\begin{align*}
  \operatorname{Intermediate\ value\ theorem}_{[0, 1]}^{0} (x^3 + x - 1) = \begin{array}{ll} \text{continuous on [0, 1]} & \text{yes} \\ \text{f(0)} & -1 \\ \text{f(1)} & 1 \\ \text{0 lies between f(0) and f(1)} & \text{yes} \\ \text{c with f(c) = 0} & \left\{-\sqrt[3]{\frac{\sqrt{93}}{18} - \frac{1}{2}} + \sqrt[3]{\frac{\sqrt{93}}{18} + \frac{1}{2}}\right\} \end{array}
\end{align*}
\section*{The full summary}

\begin{equation*}
  h(x) \coloneqq  x e^{-x}
\end{equation*}
Analyse $h$.

\end{document}
